\chapter{cna-samples and cna-examples}
\label{ch:samples-and-examples}

This book's ecosystem chapter introduced both of these sibling repositories briefly, as two
different kinds of proof that CNA works. This chapter returns to them from a practical
angle: what each one is actually useful for, if you are learning CNA or evaluating it for a
port.

\section{cna-samples: historical fidelity, and a precise measurement of one gap}

\texttt{cna-samples} is a direct C\texttt{++} port of the official Microsoft XNA Game
Studio~4.0 sample collection --- eighty-six samples in the full upstream set --- licensed
Ms-PL specifically because it is derived Microsoft content, unlike \texttt{cna-examples}
below. Building it follows the same sibling-checkout pattern as CNA itself
(Chapter~\ref{ch:building-cna}): \texttt{../cna} and \texttt{../sharp-runtime}
must sit alongside it on disk.

Its real value to a reader of this book is as a precise, checkable measurement of the
\texttt{.fx} bytecode gap from Chapter~\ref{ch:shader-fx-gap}: 23 of the 86 samples fail to
run today. Most, but not quite all, of those 23 failures trace to that same root cause --- a
compiled custom effect the content pipeline cannot load; the exact breakdown is worth stating
precisely rather than rounding up to ``all of them,'' and the next section does exactly that.
If you are trying to judge, before starting a port, whether an existing XNA game is likely to
run on CNA without hitting this gap, \texttt{cna-samples} is the closest thing to a real,
checkable answer available: a sample that runs is evidence the gap does not block content
shaped like it; a sample that fails is further, specific confirmation of exactly which content
shapes still do.

\section{The 23 failures, by actual root cause}

Checking every one of the 23 blocked samples' own \texttt{missing.md} write-ups directly,
rather than assuming they all share one cause, finds three distinct gaps, not one. Sixteen ---
\texttt{DistortionSample}, \texttt{NonPhotoRealistic}, \texttt{NormalMapping},
\texttt{PerPixelLighting}, \texttt{VertexLighting}, \texttt{ShadowMapping},
\texttt{BillboardSample}, \texttt{InstancedModel}, \texttt{ShatterEffect},
\texttt{Particles3D}, \texttt{XmlParticles}, \texttt{NetRumble}, \texttt{ShipGame},
\texttt{ColorReplacement}, \texttt{BloomSample}, and \texttt{CustomModelEffect} --- do trace to
the \texttt{.fx} bytecode gap this book has already named. \texttt{Spacewar} is blocked on a
compound of gaps that includes a custom shader among others (a custom \cnaclass{Model},
\cnaclass{RenderTarget2D} content, and XACT audio, per its own write-up), which this book counts
in the same bucket since the shader gap alone would already block it. That leaves seven
failures with a genuinely different cause. Four --- \texttt{CustomModelAnimation},
\texttt{SkinningSample}, \texttt{SkinnedModelExtensions}, and \texttt{CPUSkinning} --- are
blocked on skeletal animation playback: CNA can render a static or rigidly-jointed
\cnaclass{Model} (Chapter~\ref{ch:models-meshes}), but has no equivalent of XNA's own
\texttt{AnimationClip} / \texttt{Keyframe} / \texttt{AnimationPlayer} pipeline, and no per-vertex
bone-weight field in the \texttt{.model.json} format at all --- a real, unimplemented CNA
capability, not merely an unconverted asset. Two more --- \texttt{SplitScreen} and
\texttt{TankOnHeightmap} --- are blocked on a narrower, related gap: CNA's \texttt{.model.json}
reader builds only a single flat root bone per model today, so a rigid multi-part mesh whose
parts must move independently (a tank's turret and wheels, in both samples' case) has nowhere
to attach a per-mesh transform, even though no animation playback is involved at all. Both of
these are smaller, more tractable gaps than the shader problem --- narrow, specific, engine-level
features with a known shape --- but they are not the same gap, and a reader judging a real
game's own portability by this project's 23-of-86 figure should check which of the three
buckets their own content actually needs before assuming the \texttt{.fx} story is the whole
picture.

\section{A sample marked ``Done'' can still hide a real bug}

The 63-of-86 samples that do build and run are not, individually, a guarantee of pixel-perfect
correctness --- and this project's own history has one concrete, instructive case of finding
that out the hard way. \texttt{SimpleAnimation} (the tank-with-independently-animated-parts
sample behind the multi-bone gap the previous section just described) was marked ``Done'' once
the underlying \cnaclass{ModelBone} hierarchy fix landed. A later session, actually building and
running it rather than trusting the status line, found two real, previously-unnoticed CNA
bugs: a systematic tank-mesh winding reversal, and a \cnaclass{GraphicsDevice} default-state
depth-occlusion defect --- both fixed, and both invisible to every earlier check that stopped at
``does it compile and does the multi-bone hierarchy resolve.'' This is the same lesson
Chapter~\ref{ch:migration-guide} draws from a different angle and worth restating precisely
here: a status label of ``Done'' answers ``does this build and run,'' not ``is this pixel-correct
against the real original,'' and this discovery was serious enough that \texttt{cna-samples}'
own maintainers responded by opening a dedicated, one-task-per-sample re-audit covering all 153
samples this project has ever catalogued --- including the 67 never even attempted --- rather
than trusting any existing ``Done'' status line at face value going forward.

\subsection{Two real framework bugs a sample found that no unit test had}

Two further findings from this same porting effort are worth naming because they are not
sample-specific workarounds at all --- they are genuine CNA framework gaps that a real,
externally-authored sample happened to exercise in a way no existing unit test did.
\texttt{LensFlare}'s own occlusion-query technique deliberately draws a small test polygon with
\cnaclass{BlendState.ColorWriteChannels} set to \texttt{None}, specifically so that polygon
stays invisible while still occluding the depth/stencil test the query reads --- and on
\texttt{EASYGL}, it renders as a solid white square instead. Reading
\texttt{EasyGLGraphicsBackend.cpp} directly confirms why: there is no \texttt{glColorMask} call
anywhere in the file at all, so \cnaclass{ColorWriteChannels} is accepted by the public API and
silently never reaches the GPU on this backend --- a real, currently open gap, not something
this specific sample can work around on its own.

A second, related finding has a more encouraging ending. The same porting effort's own bug
list once read that \texttt{GraphicsDevice::Clear(Color)}'s single-argument overload --- the
overload every 3D sample in the repository actually uses --- never cleared the depth buffer at
all, only the color target, diverging from real XNA/FNA's own behavior (which clears
target/depth/stencil together through this exact overload). Reading the real, current
implementation directly shows this has since been fixed: the single-\texttt{Color} overload now
forwards to the full \texttt{ClearOptions::Target} / \texttt{DepthBuffer} / \texttt{Stencil}
form, using the device's own \emph{current} viewport \texttt{MaxDepth}
as the clear depth --- matching FNA's real \texttt{Viewport.MaxDepth} reference exactly rather
than a hardcoded \texttt{1.0}, per the fix's own explanatory comment. Both findings make the
same point about this book's own methodology as much as about CNA itself: a real, externally
authored sample exercising a code path incidentally (an occlusion-query trick, an ordinary
\texttt{Clear} call every 3D scene makes) surfaces gaps a unit test built specifically to check
one property in isolation would never think to combine.

\subsection{The largest framework bug a sample ever found: a project-wide vertex-corruption defect}

The two findings above are real, but modest in blast radius --- each affects one code path.
Porting \texttt{InverseKinematics} surfaced something categorically bigger: a silent
vertex-data corruption bug reachable by \emph{every} \cnaclass{Model} loaded through
\cnaclass{ContentManager} with a particular, extremely common vertex layout, confirmed as
the likely true cause of a symptom family this project had been chasing across several
unrelated samples for multiple sessions.

The investigation itself is worth walking through, because the isolation technique is as
instructive as the bug. A straightforward port of a 418-vertex cylinder mesh built and
loaded without error --- \cnaclass{ContentManager} reported the correct mesh count, vertex
count, and primitive count --- but the model never rendered, at any camera distance, scale,
or cull mode, even reduced to one full-scale, unlit, untextured triangle drawn through the
ordinary \cnaclass{Model} / \cnaclass{ModelMesh}\allowbreak{}\texttt{::Draw()} path. A
hand-built triangle bypassing the content pipeline entirely failed identically, and so did
an already-shipped, already-working \texttt{.model.json} loaded fresh into the same code
--- ruling out this one mesh's own data as the cause. A much larger model (3{,}252
vertices) rendered correctly through the exact same code path at every scale tested,
including scaled down to match the failing cases' size --- the detail that pointed at the
\emph{reader} itself, not any one sample's asset or draw call, and hinted that whatever was
wrong was sensitive to something about the failing meshes specifically, not to rendering in
general.

The root cause, confirmed by direct source read and a standalone \texttt{sizeof()} probe
rather than guessed: every CNA vertex struct
(\cnaclass{VertexPositionColor}, \cnaclass{VertexPositionTexture},
\cnaclass{VertexPositionColorTexture}, \cnaclass{VertexPositionNormalTexture}) publicly
inherits the polymorphic \cnaclass{IVertexType}, so every instance carries an 8-byte vtable
pointer real XNA's own clean, unpadded C\# layouts never had:

\begin{lstlisting}
sizeof(VertexPositionColor)         = 40   (XNA/clean size: 16)
sizeof(VertexPositionTexture)       = 32   (XNA/clean size: 20)
sizeof(VertexPositionColorTexture)  = 56   (XNA/clean size: 24)
sizeof(VertexPositionNormalTexture) = 40   (XNA/clean size: 32)
\end{lstlisting}

The reader, at the time, selected which typed \texttt{VertexBuffer::SetData} overload to
call for a mesh by comparing the \texttt{.model.json}'s own declared \texttt{"vertexStride"}
--- always one of the clean XNA sizes, since every offline conversion tool writes the
tightly-packed layout, never the vtable-inflated one --- directly against \texttt{sizeof()}
of these now-inflated structs. None of the four clean stride values equals its
\emph{intended} struct's real, inflated size any more --- \textbf{except one}: a declared
stride of 32 happens, by pure coincidence of \cnaclass{VertexPositionTexture}'s own
particular field-count-plus-padding arithmetic (8-byte vtable + 12-byte position + 8-byte
texture-coordinate, padded to 32 for 4-byte alignment), to equal \texttt{sizeof(Vertex}\allowbreak{}
\texttt{PositionTexture)} --- which is also 32. Every single \texttt{Content.Load<Model>}-based
sample in the entire \texttt{cna-samples} repository declares stride 32, because every
conversion tool only ever emits \cnaclass{VertexPositionNormalTexture} data. The reader
therefore always silently dispatched to the \emph{wrong} branch, reinterpreting raw,
vtable-free file bytes (laid out as position[12] + normal[12] + texcoord[8], with no
vtable at all) as an array of real, vtable-shifted \cnaclass{VertexPositionTexture} objects
--- reading each vertex's position and texture-coordinate fields from byte offsets that
matched neither the file's actual layout nor the struct's own intended one, silently
corrupting the GPU-bound vertex data for every stride-32 model in the repository, not just
the one sample that happened to surface it.

This is very likely the true root cause of a symptom family this project had already spent
several sessions attributing, without ever directly inspecting the rasterizer itself, to
``near-plane clipping'' across multiple unrelated samples: a corrupted, effectively
arbitrary byte-offset reinterpretation of a mesh's own position data produces exactly the
two symptoms already on record --- a degenerate thin line, if the corrupted positions still
land roughly in the same region of space, or full invisibility, if they land outside the
view frustum entirely --- without any actual clip-space or projection defect being involved
at all.

The fix landed in \texttt{cna} itself (confirmed by direct source read of
\texttt{ContentManager.cpp}'s \texttt{BuildVertexBufferFromRawBytes()}), and its own code
comment explicitly credits where the bug was actually found:

\begin{lstlisting}
// Task 927: `stride` is always one of XNA's own "clean" (tightly packed, no
// vtable) sizes -- 16/20/24/32/52/56 -- ... Every CNA vertex struct below
// publicly inherits the polymorphic IVertexType (a vtable pointer XNA's own
// C# interface never carried), inflating its real sizeof() past that clean
// size -- comparing `stride` against sizeof(...) directly, then
// reinterpret_cast-ing the tightly-packed bytes as an array of those
// inflated structs, silently reads every vertex field from the wrong byte
// offset (confirmed the true cause of a long-tracked "near-plane-clipping"/
// invisible-model symptom family in ../cna-samples). Read each vertex's
// fields explicitly from the buffer's own known-clean offsets instead...
\end{lstlisting}

The fix reads each vertex's fields explicitly at known-clean offsets and constructs real,
normally-initialized C\texttt{++} objects field by field --- letting the compiler resolve
member layout correctly --- rather than reinterpreting raw bytes as an inflated struct at
all; two further, larger skinned/PBR stride values sidestep the whole class of risk by
uploading raw bytes directly (\texttt{SetDataRaw()}) instead of ever constructing a
vtable-carrying struct from file data. Rebuilding the full sample suite against this fix and
re-running a previously-affected sample confirmed it live: a tank mesh previously invisible
at that sample's own camera distance rendered as a complete, correctly-shaped solid
silhouette afterward, no other change. This is, by a wide margin, the most consequential
single bug this project's sample-porting effort has ever surfaced --- not because any one
sample's code was unusual, but precisely because it was not: an ordinary, correctly-written
port of ordinary content, run far enough to notice that ``the mesh loads without error but
never appears'' is not the same claim as ``the mesh loaded correctly.''

\subsection{The 67 samples beyond the 86: not forgotten, each excluded for a named, structural reason}

The ``153 samples this project has ever catalogued'' figure above is worth breaking down
precisely rather than leaving as one round number, because the 67 beyond the official XNA
Game~Studio~4.0 collection's own 86 are not an unattempted backlog --- every single one carries
an individually recorded, structural reason it is permanently out of CNA's own scope, not a gap
awaiting engineering effort. A first group --- Xbox~LIVE Avatar animation-blending demos
(\texttt{AvatarAnimationBlending}, \texttt{AvatarMultipleAnimations},
\texttt{AvatarShadows}, and similar), each depending on the same real Xbox~LIVE Avatar
body/animation content service this book's own Avatar chapter already covers as permanently
retired --- is explicitly revisitable only if CNA's own opt-in substitute-body rendering
(Chapter~\ref{ch:avatar}) is ever judged faithful enough to stand in for it, a decision already
made and recorded, not merely deferred by omission. A second group is structural for an entirely
different reason: WinForms design-time tools (\texttt{BitmapFontMaker}, \texttt{CurveEditor},
\texttt{XnaGraphicsProfileChecker}), content-pipeline-extension-only projects with no
executable at all, and real Windows-Phone-hardware demos with no portable equivalent
(\texttt{PushNotifications}, \texttt{SavingEmbeddedImages}) are not \texttt{Game}s in any sense
this book's own API surface could run at all, regardless of CNA's own capability. A third,
larger group of 42 samples was never even given an individual tracking entry, precisely because
each one shares its own exclusion reason identically with every other sample in its group: pre-%
4.0 XNA archives (2.0/3.0/3.1-era samples this repository's own scope never covered to begin
with), pure art/rigging asset packs with no C\# code at all, a Visual Basic-language duplicate of
an already-ported C\# sample, Silverlight/Windows-Phone-native lifecycle demos with no XNA
\texttt{Game} anywhere in them, and a handful of unofficial, third-party community kits outside
the official Microsoft collection this project deliberately scopes itself to. None of these 67
carry a ``revisit trigger'' pointing back at CNA's own engineering backlog the way the 23
shader-blocked or 6 animation/multi-bone-blocked samples above do --- the honest, checkable
distinction this book keeps returning to is exactly the one between ``still open, waiting on
CNA'' and ``permanently out of scope for a reason that has nothing to do with CNA's own
capability at all,'' and this 67-sample breakdown is precisely that distinction, applied
completely.

\section{cna-examples: a coherent, in-app tour, and a project still early}

\texttt{cna-examples} answers a different question, and is deliberately original,
MIT-licensed work rather than a Microsoft-derived port --- precisely so it can carry a
license \texttt{cna-samples} cannot. Conceptually modeled on JavaFX's own \texttt{Ensemble}
sample browser, it is a single, cross-platform application (desktop, Web, and mobile) that
lets a user navigate from a top-level area (Input, Audio, Devices, Net, Media, 2D Graphics,
3D Graphics, and so on) down through categories to individual, runnable demonstrations, all
inside one binary rather than as eighty-six separate small programs.

At the time of writing, this project is an early skeleton: the application boots, shows its
home menu, and navigates from an area down to a category, but individual demo screens
themselves are not yet implemented, and only the \texttt{EASYGL} graphics backend is
currently targeted. This book states that plainly rather than describing the project by its
aspiration --- it is worth knowing about specifically because its \emph{intended} final shape
(one coherent, navigable catalog spanning the 2D/3D graphics machinery of Parts~III--IV and
the input/audio/devices/networking surface of Part~V) is a genuinely useful thing to watch for
as it matures, not because it is usable as a comprehensive demo catalog today.

\section{Reading them together}

The two repositories are complementary rather than redundant. \texttt{cna-samples} tells you
whether \emph{historical, real XNA content} runs, at the granularity of eighty-six real,
external samples CNA did not design around. \texttt{cna-examples} is intended to tell you
whether \emph{CNA's own capabilities}, area by area, are demonstrable in one coherent
application --- a question \texttt{cna-samples}, being a fixed port of pre-existing content,
cannot answer on its own, since it can only ever demonstrate what the original eighty-six
samples happened to exercise.
